% Edited conversion of source pp. 21--22.
\ReportHeading
  {Методика оптимізації режимів термообробки елементів машинобудівних конструкцій}
  {Є.~М. Ірза}
  {Інститут прикладних проблем механіки і математики ім. Я.~С. Підстригача НАН України, вул. Наукова, 3-б, Львів, 79060, Україна}
  {Evgen_Irza@ukr.net}

Елементи машинобудівних конструкцій часто піддають тепловим процесам, пов’язаним із нагріванням і подальшим охолодженням. Температурні поля спричиняють напруження, які можуть перевищувати допустимі значення, призводити до руйнування або погіршувати експлуатаційні властивості конструкції. Водночас термообробка є енергозатратною. Тому необхідні ефективні математичні методики оптимального керування режимами нагрівання--охолодження за заданими критеріями та з урахуванням обмежень на напружено-деформований стан і технологічні умови.

Розв’язування цього класу задач охоплює три основні етапи:
\begin{itemize}[nosep,leftmargin=2em]
  \item математичну постановку задачі оптимізації;
  \item розроблення чисельного алгоритму пошуку оптимального розв’язку;
  \item реалізацію алгоритму в програмному комплексі.
\end{itemize}

Математична постановка задачі оптимізації передбачає:
\begin{itemize}[nosep,leftmargin=2em]
  \item вибір параметрів стану;
  \item формулювання математичних залежностей, що описують поведінку елементів конструкцій за заданих умов термообробки;
  \item вибір функціонала та критерію оптимізації;
  \item вибір функцій керування процесом термообробки;
  \item формування математичних обмежень на параметри стану та функції керування.
\end{itemize}

Розроблення оптимального режиму термообробки складається з побудови алгоритмів розв’язування прямих задач і власне задачі оптимізації.

У запропонованому підході прямі задачі розв’язують методом зважених залишків у поєднанні з методом скінченних елементів. Алгоритм включає дискретизацію області, апроксимацію невідомих функцій на скінченному елементі, побудову системи алгебраїчних рівнянь відносно вузлових значень і розв’язування цієї системи.

Алгоритм оптимізації побудовано на принципі поетапної параметричної оптимізації. Його реалізовано в програмному комплексі, що містить:
\begin{itemize}[nosep,leftmargin=2em]
  \item базу даних геометрії області, яку займає тіло;
  \item модуль дискретизації області;
  \item бази даних граничних умов і фізичних характеристик матеріалу;
  \item розрахункові модулі для прямих та оптимізаційних задач;
  \item модулі візуалізації результатів.
\end{itemize}

Розрахункову програмну оболонку побудовано за принципами загальності, відкритості, модульності та компактності.
